% Zone „Das kennst du schon“ – aus bank/trigonometrische-funktionen/zone.jsonl (zusammenbau v0.1)
\einheitenkopf[zone][Kennst du schon]{Das kennst du schon}
\setcounter{aufgabe}{0}
%% TODO Grundvorstellungs-Aufgabe als erste Hauptnummer der Zone (2.8) fehlt in der Bank

%% TODO Titel: Ich-kann-Satz fehlt in der Bank (Platzhalter: Kettenname) – Nr. 1
%% TODO Anweisung: ein Satz über der ersten Teilaufgabe fehlt in der Bank – Nr. 1
\begin{aufgabe}{Sinus, Kosinus und Tangens als Seitenverhältnisse im rechtwinkligen Dreieck, Wert zu einem Winkel mit dem Taschenrechner im Grad-Modus}
%% TODO Darstellung für schwach fehlt in der Bank (trigonometrische-funktionen-zone-f1-v1; Abschnitt „Für schwache Schüler“)
\swz{Gegenkathete 2 cm, Hypotenuse 5 cm – wie groß ist sin $\alpha$?}{}
%% TODO Darstellung für schwach fehlt in der Bank (trigonometrische-funktionen-zone-f1-v3; Abschnitt „Für schwache Schüler“)
\swz{sin 37° mit dem Taschenrechner im Grad-Modus – auf zwei Stellen?}{}
\end{aufgabe}

%% TODO Titel: Ich-kann-Satz fehlt in der Bank (Platzhalter: Kettenname) – Nr. 2
%% TODO Anweisung: ein Satz über der ersten Teilaufgabe fehlt in der Bank – Nr. 2
\begin{aufgabe}{Taschenrechner}
\begin{teile}
\teil Winkel in Grad – welchen Modus stellst du ein?\\ \kreuz{DEG}\\ \kreuz{RAD}
\end{teile}
%% TODO Darstellung für schwach fehlt in der Bank (trigonometrische-funktionen-zone-f2-v3; Abschnitt „Für schwache Schüler“)
\swz{sin $\alpha = 0{,}42$ – wie groß ist $\alpha$? Nutze die Umkehrtaste und runde auf eine Stelle.}{}
\end{aufgabe}

%% TODO Titel: Ich-kann-Satz fehlt in der Bank (Platzhalter: Kettenname) – Nr. 3
%% TODO Anweisung: ein Satz über der ersten Teilaufgabe fehlt in der Bank – Nr. 3
\begin{aufgabe}{Koordinatensystem mit vier Quadranten, Punkte eintragen und ablesen, Achsen benennen und einteilen}
\swz{Lies die Koordinaten von A ab.}{\begin{ksys}[xmin=-4,xmax=4,ymin=-4,ymax=4,ablesen]\punkt{3}{2}{A}\end{ksys}}
\swz{Lies die Koordinaten von C ab.}{\begin{ksys}[xmin=-4,xmax=4,ymin=-4,ymax=4,xstep=0.5,ystep=0.5,ablesen]\punkt{-1.5}{-2.5}{C}\end{ksys}}
\end{aufgabe}

%% TODO Titel: Ich-kann-Satz fehlt in der Bank (Platzhalter: Kettenname) – Nr. 4
%% TODO Anweisung: ein Satz über der ersten Teilaufgabe fehlt in der Bank – Nr. 4
\begin{aufgabe}{Kreis mit Mittelpunkt und Radius zeichnen, Umfang mit Pi berechnen, Mittelpunktswinkel und Kreisbogen als Anteil des Vollkreises}
%% TODO Darstellung für schwach fehlt in der Bank (trigonometrische-funktionen-zone-f4-v1; Abschnitt „Für schwache Schüler“)
\swz{Radius 3 cm – wie lang ist der Umfang? Gib ihn mit $\pi$ und gerundet an.}{}
%% TODO Darstellung für schwach fehlt in der Bank (trigonometrische-funktionen-zone-f4-v3; Abschnitt „Für schwache Schüler“)
\swz{Radius 4 cm, Mittelpunktswinkel 45° – wie lang ist der Kreisbogen? Runde auf zwei Stellen.}{}
\end{aufgabe}

%% TODO Titel: Ich-kann-Satz fehlt in der Bank (Platzhalter: Kettenname) – Nr. 5
%% TODO Anweisung: ein Satz über der ersten Teilaufgabe fehlt in der Bank – Nr. 5
\begin{aufgabe}{Winkel messen, zeichnen und benennen, Vollwinkel, gestreckter und rechter Winkel, Winkel über neunzig Grad als stumpfe Winkel}
%% TODO Darstellung für schwach fehlt in der Bank (trigonometrische-funktionen-zone-f5-v1; Abschnitt „Für schwache Schüler“)
\swz{Wie groß ist ein gestreckter Winkel?}{}
\swa{Zeichne am Strahl einen Winkel von 135° ein.}{\winkelstrahl}
\end{aufgabe}

%% TODO Titel: Ich-kann-Satz fehlt in der Bank (Platzhalter: Kettenname) – Nr. 6
%% TODO Anweisung: ein Satz über der ersten Teilaufgabe fehlt in der Bank – Nr. 6
\begin{aufgabe}{Bruchteile eines Vollkreises als Bruch und als Dezimalzahl (halb, viertel, drittel)}
%% TODO Darstellung für schwach fehlt in der Bank (trigonometrische-funktionen-zone-f6-v1; Abschnitt „Für schwache Schüler“)
\swz{Ein Viertel des Vollkreises – wie viel Grad?}{}
%% TODO Darstellung für schwach fehlt in der Bank (trigonometrische-funktionen-zone-f6-v3; Abschnitt „Für schwache Schüler“)
\swz{72° – welcher Anteil des Vollkreises? Gib ihn als Bruch und als Dezimalzahl an.}{}
\end{aufgabe}

%% TODO Titel: Ich-kann-Satz fehlt in der Bank (Platzhalter: Kettenname) – Nr. 7
%% TODO Anweisung: ein Satz über der ersten Teilaufgabe fehlt in der Bank – Nr. 7
\begin{aufgabe}{Wertetabelle anlegen, Wertepaare als Punkte eintragen, den Graphen als durchgehende Linie zeichnen, Funktionswert und Argument ablesen, Punktprobe}
\swa{$y = 2x$ – fülle die Wertetabelle aus.}{\wertetabelle{x}{y}{0,1,2,3}}
\swa{$y = x^2 - 1$ – fülle die Wertetabelle aus.}{\wertetabelle{x}{y}{-2,-1,0,1,2}}
\end{aufgabe}

%% TODO Titel: Ich-kann-Satz fehlt in der Bank (Platzhalter: Kettenname) – Nr. 8
%% TODO Anweisung: ein Satz über der ersten Teilaufgabe fehlt in der Bank – Nr. 8
\begin{aufgabe}{Achseneinteilung wählen, wenn beide Achsen verschiedene Einheiten und Schrittweiten haben}
%% TODO Darstellung für schwach fehlt in der Bank (trigonometrische-funktionen-zone-f8-v1; Abschnitt „Für schwache Schüler“)
\swz{Rechtsachse: 0 bis 60 Minuten auf 6 cm – wie viele Minuten je Zentimeter?}{}
%% TODO Darstellung für schwach fehlt in der Bank (trigonometrische-funktionen-zone-f8-v3; Abschnitt „Für schwache Schüler“)
\swz{Rechtsachse: 0 bis 2,5 Stunden auf 10 cm – wie viel je Zentimeter?}{}
\end{aufgabe}

%% TODO Titel: Ich-kann-Satz fehlt in der Bank (Platzhalter: Kettenname) – Nr. 9
%% TODO Anweisung: ein Satz über der ersten Teilaufgabe fehlt in der Bank – Nr. 9
\begin{aufgabe}{Achsensymmetrisch und punktsymmetrisch als Wörter, Symmetrieachse und Symmetriezentrum am Graphen}
\begin{teile}
\teil $y = x^2$ – welche Symmetrie?\\ \kreuz{achsensymmetrisch zur Hochachse}\\ \kreuz{punktsymmetrisch zum Ursprung}
\end{teile}
%% TODO Darstellung für schwach fehlt in der Bank (trigonometrische-funktionen-zone-f9-v3; Abschnitt „Für schwache Schüler“)
\swz{$y = (x - 2)^2$ – wo liegt die Symmetrieachse?}{}
\end{aufgabe}

%% TODO Titel: Ich-kann-Satz fehlt in der Bank (Platzhalter: Kettenname) – Nr. 10
%% TODO Anweisung: ein Satz über der ersten Teilaufgabe fehlt in der Bank – Nr. 10
\begin{aufgabe}{Parameter in einer Funktionsgleichung erkennen und ihre Wirkung auf den Graphen beschreiben (Streckung, Stauchung, Verschiebung), Graph zu Gleichung zuordnen}
\begin{teile}
\teil Welche Parabel ist schmaler?\\ \kreuz{$y = 3x^2$}\\ \kreuz{$y = x^2$}
\end{teile}
%% TODO Darstellung für schwach fehlt in der Bank (trigonometrische-funktionen-zone-f10-v3; Abschnitt „Für schwache Schüler“)
\swz{$y = -0{,}25x^2$ – wie wirkt der Faktor vor $x^2$?}{}
\end{aufgabe}

%% TODO Titel: Ich-kann-Satz fehlt in der Bank (Platzhalter: Kettenname) – Nr. 11
%% TODO Anweisung: ein Satz über der ersten Teilaufgabe fehlt in der Bank – Nr. 11
\begin{aufgabe}{Größtwert, Kleinstwert und Spannweite einer Tabelle entnehmen, Mittelwert bilden}
%% TODO Darstellung für schwach fehlt in der Bank (trigonometrische-funktionen-zone-f11-v1; Abschnitt „Für schwache Schüler“)
\swz{Temperaturen 4, 9, 7, 11, 6 (in °C) – Größtwert und Kleinstwert?}{}
%% TODO Darstellung für schwach fehlt in der Bank (trigonometrische-funktionen-zone-f11-v3; Abschnitt „Für schwache Schüler“)
\swz{−2, 5, 3, −1, 4 (in °C) – wie groß ist der Mittelwert?}{}
\end{aufgabe}

%% TODO Titel: Ich-kann-Satz fehlt in der Bank (Platzhalter: Kettenname) – Nr. 12
%% TODO Anweisung: ein Satz über der ersten Teilaufgabe fehlt in der Bank – Nr. 12
\begin{aufgabe}{Taschenrechner – Fehler finden}
\swz[3]{Tom rechnet 8 · sin 35° und erhält −3,43. Finde den Fehler und rechne richtig.}{}
\end{aufgabe}

%% TODO Titel: Ich-kann-Satz fehlt in der Bank (Platzhalter: Kettenname) – Nr. 13
%% TODO Anweisung: ein Satz über der ersten Teilaufgabe fehlt in der Bank – Nr. 13
\begin{aufgabe}{Taschenrechner – selbst rechnen}
%% TODO Darstellung für schwach fehlt in der Bank (trigonometrische-funktionen-zone-f2-v6; Abschnitt „Für schwache Schüler“)
\swz{6 · cos 48° – auf zwei Stellen? Prüfe vorher den Modus.}{}
\end{aufgabe}
